% Laget automatisk av elevark/lag_elevark.py fra programmering/python-bevegelse-maledata.html. Endringer her blir overskrevet.
\documentclass[11pt]{article}
\usepackage{elevark}
\setlength{\svarhoyde}{1.4cm}
\begin{document}
\elevarkhode{Fysikk 1 · Rettlinjet bevegelse}{Fart og akselerasjon fra måledata}{https://lektorbo.no/programmering/python-bevegelse-maledata.html}{Koden på nett}

Fra en video eller en bevegelsessensor får vi posisjonen til en gjenstand ved ulike tidspunkter. Men hvordan finner vi farten og akselerasjonen? Her lar vi Python regne ut stigningstall mellom målepunktene, og bruker det på en sprettball.

\section*{Oppgaver}

\begin{oppgaveliste}
\oppgave \fase{test}{Kjør}Kjør steg 1. Når slippes ballen, og når treffer den gulvet første gang? Beskriv formen på de tre grafene mellom sprettene og i selve sprettene.
\svar
\oppgave \fase{forklar}{Forklar}Hvorfor er grafen for farten mer ujevn enn grafen for posisjonen, og grafen for akselerasjonen enda mer ujevn?
\svar
\oppgave \fase{test}{Regn}Kjør steg 2. Hvorfor blir middelverdien av akselerasjonen fra steg 1 dårligere enn andregradstilpasningen? Prøv tidsrommet 0,13–0,37~s i stedet.
\svar
\oppgave \fase{test}{Regn}Kjør steg 3. Hva er forholdet mellom to toppunkter som følger etter hverandre? Hva blir forholdet mellom farten rett etter og rett før et sprett?
\svar
\oppgave \fase{forklar}{Forklar}Hvor mange prosent av den mekaniske energien forsvinner i hvert sprett? Hvor blir det av energien?
\svar
\oppgave \fase{utfordring}{Utfordring}Bruk energibevaring til å regne ut farten rett før det første spretten, og høyden ballen skal nå etter spretten hvis sprettallet er 0,80. Sammenlign med det første toppunktet.
\svar
\end{oppgaveliste}

\end{document}
